\documentclass[10pt,a4paper]{amsart}
\usepackage[margin=19mm]{geometry}
\usepackage{amsmath,amssymb,mathtools}
\usepackage[scr=rsfs,cal=euler]{mathalfa}
\usepackage{xfrac}
\usepackage[unicode,hidelinks,bookmarks=false]{hyperref}
\newcommand{\ie}{i.e.\ }
\newcommand{\cf}{cf.\ }
\newcommand{\resp}{respectively\ }
\newcommand{\Subsection}{\subsection}
\providecommand{\nobreakdash}{\nobreak\hbox{-}}
\setlength{\emergencystretch}{2em}
\multlinegap0pt
\theoremstyle{plain}
\newtheorem{thm}{Theorem}
\newtheorem{prop}[thm]{Proposition}
\theoremstyle{definition}
\newtheorem{de}[thm]{Definition}
\renewcommand{\Vec}{\mathbf{Vec}}
\newcommand{\ZZ}{{\mathbb Z}}
\newcommand{\Hom}{\operatorname{Hom}\nolimits}
\title[Reduction of Noetherianity to the prime-field case]{A Tensor Restriction Theorem over Finite Fields: reduction of Noetherianity for polynomial representations to the prime-field case (Proposition 4.5.1)}
\author{Andreas Blatter, Jan Draisma, Filip Rupniewski}
\date{Source: arXiv:2211.12319v1 (CC BY 4.0)}
\begin{document}
\maketitle
\begin{abstract}
Restriction is a natural quasi-order on $d$-way tensors. We establish
a remarkable aspect of this quasi-order in the case of tensors over a
fixed finite field---namely, that it is a well-quasi-order: it admits no
infinite antichains and no infinite strictly decreasing sequences. This
result, reminiscent of the graph minor theorem, has important consequences
for an arbitrary restriction-closed tensor property $X$. For instance,
$X$ admits a characterisation by finitely many forbidden restrictions
and can be tested by looking at subtensors of a fixed size.  Our
proof involves an induction over polynomial generic representations,
establishes a generalisation of the tensor restriction theorem to other
such representations (e.g.~homogeneous polynomials of a fixed degree),
and also describes the coarse structure of any 
restriction-closed property. 
\end{abstract}


\noindent\emph{Source and licence.} A Tensor Restriction Theorem over Finite Fields. Version arXiv:2211.12319v1 (CC BY 4.0), distributed under the Creative Commons Attribution 4.0 International licence, http://creativecommons.org/licenses/by/4.0/, as recorded in the arXiv licence metadata for this exact version and shown on the abstract page https://arxiv.org/abs/2211.12319v1. Full licence notice: licenses/arxiv-2211.12319-CC-BY-4.0.txt.\par\smallskip

\noindent\textit{Editorial scope.} Counted unit: the reduction of the Noetherianity theorem for polynomial generic representations over a finite field to the case of a prime field, printed in the article as Proposition 4.5.1 (source0.tex lines 1687--1691, source label \texttt{prop:PrimeField}), delivered together with the article's complete original proof of it (source0.tex lines 1693--1715, one \texttt{proof} environment of 23 lines, adjacent to the statement, with no gap and no deferral). No mathematics is written, repaired or re-derived: statement and proof are the article's own text line for line, in the article's own notation ($K$, $F$, $e=\dim_F K$, $U_F$, $P_F$, $X_F$, $\Vec$, chains $X_1\supseteq X_2\supseteq\ldots$), and no line of the retained spans is omitted, substituted or re-flowed. Required context retained uncounted: the article's abstract (lines 41--55, reproduced verbatim as the wrapper's abstract); the article's subsection 1.5 heading with its definition of a generic representation (lines 370--377, printed as Definition 1.5.1); the article's own subsection heading and its definition of a polynomial generic representation together with its degree (lines 714 and 720--726, printed as Definition 2.2.1); the article's section and subsection headings for Noetherianity and its definition of a {\em subset} of a polynomial representation (lines 1523--1526 and 1533--1537, the definition printed as Definition 4.1.1); the article's Noetherianity theorem, which the counted statement names (lines 1539--1544, printed as Theorem 4.1.2); the article's own subsection heading for the reduction (line 1685); and the article's boxed remark recording the consequence of the proposition (lines 1717--1723). Editorial formatting only, in addition: one editorial section heading before the first retained excerpt; an AMS \texttt{amsart} preamble that restates the macros and theorem environments the retained lines use (the article's \texttt{\textbackslash Vec} from its preamble, and \texttt{\textbackslash ZZ}, \texttt{\textbackslash Hom} and the environments \texttt{thm}, \texttt{prop}, \texttt{de} as the article's own style file \texttt{all2021.sty} declares them) -- that style file is neither shipped nor input, and the \texttt{de} environments therefore lose only the decorative diamond end-of-environment symbol that the style file attaches to them; and sequential numbering of the retained environments, so that the four retained definition-like environments and the retained theorem and proposition are numbered 1, 2, 3, 4 and 5 in order of appearance, whereas the article prints them as Definition 1.5.1, Definition 2.2.1, Definition 4.1.1, Theorem 4.1.2 and Proposition 4.5.1; the cross-references inside the retained text resolve to the wrapper's numbering. No citation and no figure occurs in any retained line, so the wrapper carries no bibliography, no bibliography entry and no asset.

\section{Generic and polynomial representations}
\subsection{Generic representations}

Let $\Vec$ be the category of finite-dimensional vector spaces over the
finite field $K$.

\begin{de}
A {\em generic representation} is a functor $F:\Vec \to \Vec$.  
\end{de}

\subsection{Polynomial generic representations over $K$}

\begin{de} \label{de:PolRep}
A generic representation $P:\Vec \to \Vec$ is called {\em polynomial}
if there exists a $d$ such that for all $U,V \in \Vec$ the map
$P:\Hom(U,V) \to \Hom(P(U),P(V))$ lies in $K[\Hom(U,V)]_{\leq d} \otimes
\Hom(P(U),P(V))$. The minimal such $d \in \ZZ_{\geq -1}$ is called the
{\em degree} of $P$ and denoted $\deg(P)$.
\end{de}

\section{Noetherianity for polynomial representations}
\label{sec:Noetherianity}

\subsection{The main result}

\begin{de} \label{de:Subset}
Let $P$ be a polynomial representation. A {\em subset} of $P$ is the data of a
subset $X(V)$ of $P(V)$ for every $V \in \Vec$, subject to the condition
that for all $V,W,\phi \in \Hom_\Vec(V,W)$, $P(\phi)X(V) \subseteq X(W)$.
\end{de}

\begin{thm}[Noetherianity] \label{thm:Noetherianity}
Let $P$ be a polynomial representation over the finite field $K$. Then
any descending chain
\[ P \supseteq X_1 \supseteq X_2 \supseteq \ldots \]
of subsets stabilises.
\end{thm}

\subsection{Reduction to the prime field case}

\begin{prop} \label{prop:PrimeField}
Suppose that Theorem~\ref{thm:Noetherianity} holds when $K$
is a prime field. Then it also holds when $K$ is an
arbitrary finite field.
\end{prop}

\begin{proof}
Let $F$ be the prime field of $K$ and set $e:=\dim_F K$. For
an $n$-dimensional $K$-vector space $U$, we write $U_F$ for the
$e \cdot n$-dimensional $F$-vector space obtained by restricting the scalar
multiplication on $U$ from $K \times U \to U$ to $F \times U \to U$.

Now let $P$ be a polynomial representation over $K$. Define a generic
representation $P_F$ over $F$ by setting, for a finite-dimensional
$F$-vector space $U$, $P_F(U):=(P(K \otimes_F U))_F$, and sending an
$F$-linear map $\phi:U \to V$ to the map $P_F(\phi)$, which is $K$-linear
and therefore also $F$-linear. It is easy to see from the definitions
that $P_F$ is polynomial of the same degree as $P$.

For a subset $X$ of $P$, we define a subset of $P_F$ via
$X_F(U):=X(K \otimes_F U)$. If $X_1 \supseteq X_2 \supseteq \ldots$ is a
chain of subsets in $P$, then $(X_1)_F \supseteq (X_2)_F \supseteq \ldots$
is a chain of subsets in $P_F$. By assumption, the latter
stabilises, say at $(X_{n_0})_F$. Then it follows that, for any
$n \geq n_0$ and any $m$, 
\[ X_{n}(K^m)=X_{n}(K \otimes_F F^m)=(X_{n})_F(F^m)
=(X_{n_0})_F(F^m)=X_{n_0}(K^m), \]
and this suffices to conclude that $X_n=X_{n_0}$. 
\end{proof}

\begin{center}
\fbox{\parbox{.8\textwidth}{
{\bf In view of Proposition~\ref{prop:PrimeField}, from
now on
we assume
that $K$ is a prime field.}}} 
\end{center}

\end{document}
